\documentclass[11pt, a4paper]{article}

\usepackage[a4paper, margin=2.5cm]{geometry}
\usepackage{fontspec}
\usepackage[portuguese, bidi=basic, provide=*]{babel}
\babelprovide[import, onchar=ids fonts]{portuguese}
\babelfont{rm}{Noto Sans}

\usepackage{booktabs}
\usepackage{enumitem}
\setlist[itemize]{label=-}
\setlist[enumerate]{leftmargin=*}
\usepackage{hyperref}
\hypersetup{
    colorlinks=true,
    linkcolor=blue,
    urlcolor=blue,
    pdftitle={Relatório da Atividade 1: Definição e Organização do Projeto Finance.io}
}

\title{\textbf{RELATÓRIO DA ATIVIDADE 1}\\Definição e Organização do Projeto de Software}
\author{
    Antonio Nilton Gomes Damasceno \and
    Fábio Ribeiro Oliveira Filho \and
    Henry Phietro Aguiar da Silva
}
\date{\today}

\begin{document}

\maketitle

\section{Integrantes da Equipe}
\begin{enumerate}
    \item Antonio Nilton Gomes Damasceno
    \item Fábio Ribeiro Oliveira Filho
    \item Henry Phietro Aguiar da Silva
\end{enumerate}

\section{Identificação do Sistema}
\textbf{Nome do Sistema:} Finance.io

\section{Descrição do Sistema}
O \textbf{Finance.io} é uma aplicação simples de gestão financeira pessoal. Seu objetivo é permitir que uma pessoa registre manualmente suas entradas e saídas de dinheiro, organize as movimentações por conta e categoria e acompanhe a situação financeira em um painel consolidado.

O produto não movimenta dinheiro de verdade. Ele funciona como uma ferramenta de registro, planejamento e consulta. O sistema também não se conecta automaticamente a bancos no MVP; os dados financeiros são informados e mantidos pelo próprio usuário.

\section{Organização e Arquitetura do Sistema}

\begin{table}[h]
\centering
\caption{Componentes e Tecnologias Utilizadas}
\vspace{0.2cm}
\begin{tabular}{lp{10cm}}
\toprule
\textbf{Componentes} & \textbf{Descrição e Tecnologias Utilizadas} \\
\midrule
Front-end & HTML, CSS, JavaScript, Vite e Tailwind. \\
Back-end & Python 3.14. \\
Banco de Dados & PostgreSQL. \\
\bottomrule
\end{tabular}
\end{table}

\section{Requisitos Principais do Sistema}

\begin{table}[h]
\centering
\caption{Requisitos Funcionais do Sistema}
\vspace{0.2cm}
\begin{tabular}{lll}
\toprule
\textbf{Identificador} & \textbf{Tipo} & \textbf{Descrição do Requisito} \\
\midrule
RF01 & Funcional & Registrar Entradas Financeiras. \\
RF02 & Funcional & Registrar Saídas Financeiras. \\
RF03 & Funcional & Edição de Registros. \\
RF04 & Funcional & Organizar as movimentações. \\
RF05 & Funcional & Acompanhar os Registros. \\
RF06 & Funcional & Oferecer Visão Geral das Finanças. \\
\bottomrule
\end{tabular}
\end{table}

\section{Links e Imagens, GitHub e Taiga}

\begin{itemize}
    \item \textbf{Repositório no GitHub:} \url{https://github.com/EduPhietro/Finance.io}
    \item \textbf{Projeto no Taiga:} \url{https://tree.taiga.io/project/eduphietro-erpa/timeline}
\end{itemize}

\end{document}